\section{U.S. Data Center Tax Incentives: Policy Landscape and Recent Changes}

Data center tax incentives in the United States are primarily established by state and local governments. In the context of crypto mining companies converting or building data centers, state-level tax incentives directly affect two major areas:

First, construction materials and equipment procurement. During a data center conversion, a mining company may need to purchase substantial quantities of tangible personal property, including servers, electrical equipment and networking equipment, to support the new business activities. These purchases may trigger sales and use tax, while some states provide exemptions if specified conditions are met.

Second, ongoing data center operations. Once a mining company's data center is completed and placed in service, land, buildings and equipment may remain subject to property tax, while high levels of electricity consumption may also be subject to sales tax, utility tax or electricity consumption tax. Accordingly, some state and local governments provide property tax relief, electricity-related tax exemptions or other economic incentives.

For many years, numerous U.S. states used tax incentives to attract data center investment. By 2026, however, public attitudes toward data center infrastructure began to shift from support toward opposition, and the policy environment changed markedly.

Sales and use tax incentives in major jurisdictions such as Texas, Virginia and Georgia remain available for qualifying projects. At the same time, as AI data centers expand rapidly, concerns over electricity prices, water resources and burdens on local residents have drawn broad attention, prompting several states to suspend data center incentives or consider similar measures. Beginning in summer 2026, Arizona imposed a three-year moratorium on new computer data center tax-relief applications, Washington eliminated sales and use tax benefits for replacement server equipment and refurbishment, and Louisiana raised the conditions attached to data center incentives, including requirements for data center companies to bear the full incremental power needs associated with their projects.

This means that crypto mining companies may still be able to access these tax incentives during their transition, but the availability of benefits is becoming more uncertain. Project location, whether a new exemption certification is required, the timing of the application, and new factors introduced by policy changes can all affect the actual tax cost of the transition.

% Transcribed from assets/figure-1.png; wording and values retained.
\begingroup
\ReportTableSetup
\begin{longtable}{|L{0.15\linewidth}|L{0.23\linewidth}|L{\dimexpr0.62\linewidth-6\tabcolsep-4\arrayrulewidth\relax}|}
\caption{Status of Data Center Tax Incentives by State}\label{tab:state-incentives}\\
\hline
\textbf{State} & \textbf{Current Status} & \textbf{Key Development} \\ \hline
\endfirsthead
\caption[]{Status of Data Center Tax Incentives by State (continued)}\\
\hline
\textbf{State} & \textbf{Current Status} & \textbf{Key Development} \\ \hline
\endhead
Arizona & New Applications Suspended & From July 1, 2026 through June 30, 2029, the Arizona Commerce Authority may not accept applications for new computer data centers. \\ \hline
Ohio & New Applications Suspended & The Ohio Tax Credit Authority has suspended consideration of new requests for data center tax exemptions. \\ \hline
Illinois & New Agreements Suspended & Beginning July 1, 2026, the state stopped processing new applications for the data center incentive program. \\ \hline
Massachusetts & New Applications Suspended & Beginning June 25, 2026, the state paused acceptance of applications for the Qualified Data Center Sales and Use Tax Exemption. \\ \hline
Washington & Scope of Exemption Narrowed & Beginning July 1, 2026, replacement server equipment no longer qualifies for the prior exemption; no new rural refurbishment certificates are issued, and urban refurbishment projects no longer qualify. \\ \hline
Virginia & Exemption Available; New Electricity Tax & The sales and use tax exemption remains available; from July 1, 2026 through June 30, 2028, a Data Center Electricity Consumption Tax of \$0.011 per kilowatt-hour applies. \\ \hline
Louisiana & Exemption Available; Conditions Pending & A June 2026 executive order requires a ratepayer and community protection framework for data center sales tax incentives, with detailed criteria under development by the relevant authorities. \\ \hline
Texas & Exemption Available; Higher Risk of Policy Changes & The state has required a comprehensive review of data center interconnection projects and has proposed working with the Legislature to repeal outdated or unnecessary data center incentives. \\ \hline

\end{longtable}
\endgroup


\FloatBarrier
